\documentclass{article}
\usepackage{amsmath,amssymb,amsthm}
\newtheorem{theorem}{Theorem}
\begin{document}
\title{Two further 683-edge profiles excluded by simultaneous direction parity}
\author{Near-extremal campaign, 2026-11-15}
\maketitle

\section{D01: Original-graph input and actual-face import}
Let $G\subseteq Q_8$ be C4-free with exactly683 edges, with either global
parity $P$ specified and the other unrestricted. The preceding proof
q8\_equality\_profiles/main.tex gives
\[
 C(P)+\Phi\le1149,\quad
 C(P)=735n_0+534n_1+333n_2+177n_3+66n_4+105n_7+345n_8,
\]
where $\Phi\ge0$ is a sum of actual $Q_6$-face losses, zero exactly when
every face has zero or two specified-side roots of local degree at most
two. We use its full proof and declared317/321 imports; this is not a
new minimum-degree assumption on an arbitrary683 graph.

\begin{theorem}
Neither of the following specified-side profiles $(n_0,\ldots,n_8)$
occurs in such a graph:
\[
 Z_3=(0,0,0,3,3,74,44,4,0),\qquad
 Z_5=(0,0,0,5,4,62,57,0,0).
\]
\end{theorem}
Both profiles have128 vertices, degree sum683 and $C(P)=1149$.
Thus every actual $Q_6$ face has low count zero or two.

\section{D02: One simultaneous mask equation}
Only the degree-three and degree-four global roots can be locally low.
Let $R(v)$ be the retained-coordinate set at such a root.
For an omitted coordinate pair $I=\{a,b\}$, a degree-three root is low
exactly when $I\cap R(v)\ne\varnothing$; a degree-four root is low exactly
when $I\subseteq R(v)$. Summing the even low counts in the four actual
fixed-bit faces for this same $I$ proves
\[
 \sum_{d(v)=3} f(R(v)) +\sum_{d(v)=4} k(R(v))=0
       \quad\hbox{in }\mathbb F_2^{28},                       \tag{1}
\]
where $f(R)_I=\mathbf1_{I\cap R\ne\varnothing}$ and
$k(R)_I=\mathbf1_{I\subseteq R}$. No root positions or fixed bits have
been normalized independently.
For a fixed coordinate $a$, the number of its seven incident pairs on
which $f(R)$ is one is seven if $a\in R$ and three otherwise. Both are
odd. For $k(R)$ with $|R|=4$, that number is three if $a\in R$ and zero
otherwise. Since both profiles have an odd number of degree-three
roots, (1) forces each coordinate to belong to an odd number of the
degree-four retained masks.

\section{D03: Three degree-three and three degree-four roots}
There are three four-masks, with total coordinate multiplicity12. Each
coordinate occurs once or three times, so exactly two coordinates occur
three times and the other six occur once. Denote the common pair by $A$;
each four-mask contains $A$ and one of three disjoint private pairs
$B,C,D$. A single coordinate permutation, applied to the entire graph,
therefore sends the three masks to15,51,195 (bits0..7), with
$A=\{0,1\},B=\{2,3\},C=\{4,5\},D=\{6,7\}$.

The complete finite check in D05 shows that the only unordered triples
of three-masks, with repetitions permitted initially, satisfying (1)
are
\[
 (13,49,193),\quad(14,50,194).
\]
They are precisely $\{a\}\cup B,\{a\}\cup C,\{a\}\cup D$ for one
common $a\in A$. Name their actual roots $x_B,x_C,x_D$.
For an omitted pair with one coordinate in $B$ and one in $C$, exactly
$x_B,x_C$ are low: neither the third three-mask nor any four-mask
contributes. Since both actual faces must have even low count, the two
roots agree on those two fixed coordinates. Varying the pair proves
they agree on all of $B\cup C$. Using $B,D$ and $C,D$ in the same way
proves that all three roots agree on every coordinate in $B\cup C\cup D$.
They can differ only on the two coordinates of $A$. But they belong
to the same global parity; after the other six bits are fixed there
are only two possible two-bit words of that parity. Three distinct
roots cannot occupy those two positions. This excludes $Z_3$.

\section{D04: Five degree-three and four degree-four roots}
There are four four-masks, with total coordinate multiplicity16. The
odd multiplicities are one or three, so exactly four coordinates have
each type. Index the four masks by $j=0,1,2,3$. Let $u_j$ count the
coordinates present only in mask$j$, and $t_j$ those absent only from
mask$j$. Then $\sum u_j=\sum t_j=4$ and mask$j$ has size
$u_j+(4-t_j)=4$, whence $u_j=t_j$.
Up to renaming these four roots and one simultaneous coordinate
permutation, the multiplicities $u_j$ are exactly
\[
 (4,0,0,0),(3,1,0,0),(2,2,0,0),(2,1,1,0),(1,1,1,1).
\]
This permits repeated retained masks, which occur in some cases.
The representative four-masks are respectively
\[
\begin{array}{rrrr}
15&240&240&240\\135&120&184&184\\195&60&204&204\\
163&156&108&172\\169&166&154&106.
\end{array}
\]
In each of these five complete normalized cases the finite check finds
no multiset of five three-masks satisfying (1). Hence no actual root
positions could satisfy the necessary equation. This excludes $Z_5$.

\section{D05: Exact finite coverage and replay}
There are exactly $\binom83=56$ possible retained three-masks. The pair
coordinates are the28 unordered coordinate pairs in lexicographic order;
encode each vector as its28-bit nonnegative integer. For D03 enumerate
all $\binom{58}3=30856$ multisets of three masks and evaluate their XOR
against $k(15)+k(51)+k(195)$; the two displayed triples are the entire
survivor list.
For D04 enumerate all $\binom{60}5=5461512$ multisets of five masks in
each of the five cases; all fail. A separately coded replay constructs
the complete set of pair XORs and the complete set of triple XORs, and
checks that the former is disjoint from target plus the latter. Every
five-element multiset splits into a pair and a triple, so these empty
intersections independently cover the full finite domain, including
repetitions. The pair/triple code-set sizes are1541 and23926.

The deterministic scripts, direct C++ source, exact input definitions,
counts and hashes are saved alongside this proof. They make no solver
UNSAT inference and do not enumerate full Q8 graphs. The analytic
normalizations prove that every original profile maps to one of the
finite inputs, and D03 retains the actual root fixed bits after the
direction check. This is a local finite-assisted proof with its earlier
imports, not the unrestricted682 bound.
\end{document}
